\documentclass[11pt]{article}

\usepackage[T1]{fontenc}
\usepackage[utf8]{inputenc}
\usepackage{amsmath,amssymb,amsthm}
\usepackage[margin=1in]{geometry}
\usepackage{microtype}
\usepackage[hidelinks]{hyperref}
\hypersetup{
  pdftitle={Restricted-Sumset Saturation with Two Missing Sums in Prime Fields},
  pdfauthor={Xiaojian Zeng}
}

\newtheorem{theorem}{Theorem}
\newtheorem{lemma}{Lemma}
\newcommand{\Fp}{\mathbb F_p}
\newcommand{\rsum}{\mathbin{\dot{+}}}

\title{Restricted-Sumset Saturation with Two Missing Sums in Prime Fields}
\author{Xiaojian Zeng\\
  \small PBC School of Finance, Tsinghua University, 100190\\
  \small \texttt{xiaojianzeng@foxmail.com}}
\date{}

\begin{document}
\maketitle

\begin{abstract}
For subsets $A,B$ of a prime field, let $A\rsum B$ consist of the sums
$x+y$ with $x\in A$, $y\in B$, and $x\ne y$.  For every prime $p\geq
11$ and $3\leq \ell\leq (p-3)/2$, we classify, up to simultaneous
affine transformations, all ordered pairs satisfying
$|B|=\ell$, $|A|=p-\ell$, and $|A\rsum B|=p-2$.  After the two missing
sums are normalized to $\{0,2\}$, every pair is indexed by a decomposition
$\ell=a+b$ into nonnegative integers.  Two such representatives are
equivalent exactly when their parameters agree up to swapping, giving
$\lfloor\ell/2\rfloor+1$ affine classes.  The main ingredient is an exact
endpoint-corrected count for cyclic step-two runs; a direct converse handles
both degenerate endpoint families.
\end{abstract}

\section{Context and statement}

The symmetric restricted-sum problem in $\mathbb Z/p\mathbb Z$ was settled
by Dias da Silva and Hamidoune~\cite{DiasDaSilvaHamidoune1994}.  For two
nonempty sets of different sizes, Alon, Nathanson, and Ruzsa proved the sharp lower
bound
\[
 |A\rsum B|\geq \min\{p,|A|+|B|-2\}
\]
and raised the corresponding inverse problem~\cite{AlonNathansonRuzsa1996}.
K\'arolyi proved that, for $|A|\geq5$ and $2|A|-3<p$, equality
$|A\rsum A|=2|A|-3$ holds exactly for arithmetic
progressions~\cite{Karolyi2005}.  More recently, Liu
and Qian classified the asymmetric critical pairs when at least one set is
an arithmetic progression~\cite{LiuQian2026}.  Li and Liang exhibited a
boundary counterexample to the ensuing general conjecture and proved a
refined theorem in the range $|A|+|B|\leq p-1$~\cite{LiLiang2026}.

The present note treats the complementary boundary $|A|+|B|=p$ when
exactly two sums are missing and $|B|$ is in the range specified below,
without assuming that either set is an arithmetic progression.  Thus its
hypotheses lie outside Li--Liang's refined range, while its conclusion is not
restricted to the arithmetic-progression case of
Liu--Qian~\cite{LiuQian2026,LiLiang2026}.  This scope comparison is not a claim of
priority, and our elementary proof does not invoke either inverse theorem.
A separate finite-Abelian-group edge-sum problem uses coset and fiber
constructions for rainbow trees~\cite{Zeng2026Rainbow}.  It concerns
pairwise-distinct edge sums rather than restricted set addition, and no result
from it is used here.

All sets below lie in $\Fp$, where $p$ is prime.  For $A,B\subseteq\Fp$,
define the asymmetric restricted sumset
\[
 A\rsum B=\{x+y:x\in A,\ y\in B,\ x\ne y\}.
\]
The pair $(A,B)$ is ordered.  Simultaneous affine equivalence means
\[
 (A,B)\longmapsto (uA+v,uB+v),\qquad u\in\Fp^\times,\quad v\in\Fp.
\]
For nonnegative integers $a,b$ with $a+b=\ell$, put
\begin{align*}
 B_{a,b}&=\{-2j:0\leq j<a\}\cup\{1+2j:0\leq j<b\},\\
 C_{a,b}&=\{2+2j:0\leq j<a\}\cup\{-1-2j:0\leq j<b\},\\
 A_{a,b}&=\Fp\setminus C_{a,b}.
\end{align*}
An index interval is empty when its upper endpoint is below its lower
endpoint, and all displayed integers are reduced modulo $p$.

\begin{theorem}\label{thm:classification}
Let $p\geq11$ be prime and let $3\leq\ell\leq(p-3)/2$.  If
$A,B\subseteq\Fp$ satisfy
\[
 |B|=\ell,\qquad |A|=p-\ell,\qquad |A\rsum B|=p-2,
\]
then $(A,B)$ is simultaneously affinely equivalent to
$(A_{a,b},B_{a,b})$ for some $a,b\geq0$ with $a+b=\ell$.  Conversely,
every displayed pair has these cardinalities and
\[
 A_{a,b}\rsum B_{a,b}=\Fp\setminus\{0,2\}.
\]
Moreover, two displayed representatives are affinely equivalent if and only
if their parameters agree up to $(a,b)\leftrightarrow(b,a)$.  Hence there
are exactly $\lfloor\ell/2\rfloor+1$ affine classes.
\end{theorem}

\section{The endpoint-corrected run lemma}

The additive translation $x\mapsto x+2$ is one directed cycle on $\Fp$.
A cyclic $+2$-run in a proper subset $X$ means a maximal consecutive block
in this cycle.

\begin{lemma}\label{lem:runs}
Let $\varnothing\ne X\subsetneq\Fp$, $|X|=\ell$, and
\[
 V=(X\setminus\{0\})\cup((X+2)\setminus\{1\}).
\]
If $|V|\leq\ell$ and $2\ell\leq p-3$, then there are unique nonnegative
integers $a,b$ with $a+b=\ell$ such that
\begin{equation}\label{eq:Xform}
 X=\{2j:0\leq j<a\}\cup\{-1-2j:0\leq j<b\}.
\end{equation}
For these parameters,
\begin{equation}\label{eq:Vform}
 V=\{2j:1\leq j\leq a\}\cup\{-1-2j:0\leq j<b\},
 \qquad |V|=\ell.
\end{equation}
\end{lemma}

\begin{proof}
Let $r$ be the number of maximal cyclic $+2$-runs of $X$.  Each run
contributes one new endpoint to $X+2$, so
\begin{equation}\label{eq:unioncount}
 |X\cup(X+2)|=\ell+r.
\end{equation}
Let $\delta_0$ indicate that $0\in X$ and $-2\notin X$, equivalently that
$0$ is a run start.  Removing $0$ from the first copy of $X$ removes it
from the union in \eqref{eq:unioncount} exactly in this event.  Similarly,
let $\delta_1$ indicate that $-1\in X$ and $1\notin X$, equivalently that
$-1$ is a run end.  Removing $1$ from $X+2$ removes it from the union
exactly in this second event.  As $0$ and $1$ are distinct,
\begin{equation}\label{eq:endpointcount}
 |V|=\ell+r-\delta_0-\delta_1.
\end{equation}
Thus $r\leq\delta_0+\delta_1\leq2$.

If $r=1$, the run starts at $0$ or ends at $-1$.  These alternatives give,
respectively,
\[
 \{0,2,\ldots,2(\ell-1)\}
 \quad\hbox{or}\quad
 \{-1,-3,\ldots,-1-2(\ell-1)\}.
\]
Both marked endpoint events cannot occur on one run: a run of length
$m\leq\ell$ starting at $0$ and ending at $-1$ would imply
$2(m-1)\equiv-1\pmod p$, whereas
\[
 1\leq2m-1\leq2\ell-1\leq p-4.
\]

If $r=2$, equation~\eqref{eq:endpointcount} forces
$\delta_0=\delta_1=1$.  The two marked endpoints lie on different runs by
the same congruence argument.  Calling the lengths of the run starting at
$0$ and the run ending at $-1$ respectively $a$ and $b$ gives
\eqref{eq:Xform} and $a+b=\ell$.  The one-run cases are included by
allowing $a=0$ or $b=0$.

Shifting \eqref{eq:Xform} and making the two deletions gives
\eqref{eq:Vform}.  Its two displayed pieces cannot meet: if they did, then
$p$ would divide $2(i+j)+1$ for some $1\leq i\leq a$ and $0\leq j<b$,
while
\[
 3\leq2(i+j)+1\leq2\ell-1\leq p-4.
\]
Each piece is internally injective as well.  Therefore $|V|=a+b=\ell$.
The maximal cyclic-run decomposition also recovers $a$ and $b$ uniquely.
These inequalities supply the needed no-wrap argument in the field.
\end{proof}

\section{Classification and converse}

\begin{proof}[Proof of Theorem~\ref{thm:classification}]
Suppose first that $(A,B)$ satisfies the hypotheses.  Its two missing sums
are distinct residues $s,t$.  Under $z\mapsto uz+v$, restricted sums
transform by $w\mapsto uw+2v$, since the affine map preserves inequality
of summands.  Choose
\[
 u=\frac{2}{t-s},\qquad 2v=-us.
\]
This normalizes the unordered missing set to $\{0,2\}$.

Work in this normalization and put $C=\Fp\setminus A$.  If
$y\in B\setminus\{0\}$, then $-y\ne y$ because $p$ is odd.  Missingness
of $0$ therefore forces $-y\in C$.  Likewise, if
$y\in B\setminus\{1\}$, missingness of $2$ forces $2-y\in C$.  Hence
\begin{equation}\label{eq:forced}
 U(B):=((-B)\setminus\{0\})\cup((2-B)\setminus\{1\})\subseteq C.
\end{equation}
As $|C|=\ell$, we have $|U(B)|\leq\ell$.  Apply
Lemma~\ref{lem:runs} to $X=-B$.  It gives unique $a,b\geq0$ with
$a+b=\ell$ and
\[
 B=B_{a,b},\qquad U(B)=C_{a,b},\qquad |U(B)|=\ell.
\]
The inclusion \eqref{eq:forced} and equality of cardinalities now give
$C=C_{a,b}$, hence $A=A_{a,b}$.

We next prove the converse without relying on cardinality alone.  A collision
between the two pieces of $B_{a,b}$ would make $p$ divide
$2(i+j)+1\leq2\ell-3$; a collision between the pieces of $C_{a,b}$ would
make $p$ divide an odd integer at most $2\ell-1$.  Internal collisions are
also excluded because $2$ is invertible and all index differences have
absolute value below $p$.  Thus
\[
 |B_{a,b}|=|C_{a,b}|=\ell,
 \qquad |A_{a,b}|=p-\ell.
\]
Equation~\eqref{eq:Vform} says $U(B_{a,b})=C_{a,b}$.  Consequently an
allowed representation of $0$ or $2$ is impossible: except for the
equal-summand choices $0+0$ and $1+1$, its $A$-summand would belong to
$U(B_{a,b})=C_{a,b}$.

It remains to represent every other residue.  Assume first that $a,b>0$.
Then $0,1\in B_{a,b}$.  Every
$w\notin C_{a,b}\cup\{0\}$ has the allowed representation $w=w+0$.
For the remaining elements of $C_{a,b}\setminus\{2\}$, use
\begin{align*}
 2i&=(2i-1)+1 &&(2\leq i\leq a),\\
 -1-2j&=(-2j-2)+1 &&(0\leq j<b).
\end{align*}
In both lines the first summand lies outside $C_{a,b}$ and differs from
$1$.  To verify this in $\Fp$, every possible collision with either piece
of $C_{a,b}$ reduces either to a nonzero odd integer of absolute value less
than $p$, or to a positive even integer at most $2\ell<p$.  Also
$-2j-2\not\equiv1$, since $3\leq2j+3\leq2\ell-1<p$.  Thus no integer-order
argument has been imported without a modular bound.

If $a=\ell$ and $b=0$, residues outside
$C_{\ell,0}\cup\{0\}$ are again represented using $0$.  For
$2\leq i\leq\ell$, set $j=\ell-i+1$ and write
\[
 2i=2(\ell+1)+(-2j).
\]
Here $-2j\in B_{\ell,0}$, while
$2\ell+2\leq p-1$ places the first summand outside $C_{\ell,0}$.  The two
summands are not congruent: their positive even difference is below $2p$,
so it is neither the odd number $p$ nor $2p$.  The affine involution
$z\mapsto1-z$ carries this endpoint family to $a=0,b=\ell$ and transforms
sums by $w\mapsto2-w$.  This proves
$A_{a,b}\rsum B_{a,b}=\Fp\setminus\{0,2\}$ in every case.

Finally, suppose $z\mapsto uz+v$ sends one displayed ordered pair to
another.  It must preserve their missing-sum set, so
\[
 \{2v,2u+2v\}=\{0,2\}.
\]
Thus either $(u,v)=(1,0)$ or $(u,v)=(-1,1)$.  The identity preserves the
unique run lengths, whereas direct calculation gives
\[
 1-B_{a,b}=B_{b,a},\qquad 1-C_{a,b}=C_{b,a}.
\]
The parameters therefore agree exactly up to swapping.  The $\ell+1$
ordered decompositions of $\ell$ have
$\lfloor\ell/2\rfloor+1$ orbits under this swap.
\end{proof}

\section*{Computational and AI assistance}

An exact Python enumeration was used only as a falsification check: it
tested every legal value of $\ell$ for $p=11,13,17,19$ after normalization.
The computation is not used in the universal proof.  OpenAI Codex assisted
with organizing the cyclic-run argument, checking endpoint and modular
collision cases, locating bibliographic records, and editing the manuscript.
The mathematical claims remain the author's responsibility; an AI-assisted
review is not peer review.

\bibliographystyle{plain}
\bibliography{references}

\end{document}
